\documentclass[11pt, letterpaper]{article}

% --- Packages ---
\usepackage{amsmath, graphicx, amsfonts, enumitem, xcolor, listings, tcolorbox}
\usepackage[left=1in, right=1in, bottom=1in, top=3cm, headheight=2cm, headsep=0.5cm, footskip=1cm]{geometry}
\usepackage{fancyhdr}

% --- Code Listing Formatting ---
\definecolor{codegreen}{rgb}{0,0.6,0}
\definecolor{codegray}{rgb}{0.5,0.5,0.5}
\definecolor{codepurple}{rgb}{0.58,0,0.82}
\definecolor{backcolour}{rgb}{0.95,0.95,0.92}

\lstset{ 
    backgroundcolor=\color{backcolour},   
    commentstyle=\color{codegreen},
    keywordstyle=\color{blue},
    numberstyle=\tiny\color{codegray},
    stringstyle=\color{codepurple},
    basicstyle=\ttfamily\footnotesize,
    breakatwhitespace=false,         
    breaklines=true,                 
    captionpos=b,                    
    keepspaces=true,                 
    numbers=left,                    
    numbersep=5pt,                  
    showspaces=false,                
    showstringspaces=false,
    showtabs=false,                  
    tabsize=2,
    language=C++
}

% --- Header Configuration ---
\pagestyle{fancy} 
\fancyhf{} 
\renewcommand{\headrulewidth}{0.4pt} 
\lhead{\raisebox{-0.5cm}{\textbf{Universidad de Antioquia}}}
\rhead{\large \textbf{Tech Brief: Multiple Shooting Implementation} \\ \normalsize Spaceflight Dynamics 2026-1} 
\cfoot{\thepage} 

\begin{document}

\begin{center}
    {\Large \textbf{Implementing the Voyager Trajectory in GMAT}} \\
    \vspace{0.2cm}
    {\large \textit{Comparing Differential Correctors vs. Yukon Optimizers}} \\
    \vspace{0.4cm}
    \textit{Prof. Juan | Aerospace Engineering Program}
\end{center}
\hrule
\vspace{0.5cm}

\section*{Laboratory Overview}
In our analytical Python analysis, we discovered the 1976--1980 Earth-Jupiter-Saturn launch window. Python calculated idealized velocity vectors assuming instantaneous gravity. In this lab, we will inject these vectors into GMAT's high-fidelity N-body numerical integrator using a Multiple Shooting method.

We will execute this in two phases to prove why Optimizers are necessary for interplanetary design:
\begin{enumerate}[noitemsep]
    \item \textbf{Phase 1 (The Targeter):} We use a Differential Corrector (DC) to simply force the two trajectory legs to meet at Jupiter. 
    \item \textbf{Phase 2 (The Optimizer):} We swap the DC for the Yukon Optimizer to stitch the trajectories together while mathematically minimizing the required flyby $\Delta V$.
\end{enumerate}

% ==============================================================================
\section*{Step 1: Global Resource Configuration}
First, we must define the coordinate system, the spacecraft, and the solver tools. Note the specific use of the \texttt{SolarSystemBarycenter}, which is required to align GMAT's reference frame with Python's SPICE outputs.

\begin{lstlisting}[caption=GMAT Resource Initialization]
% 1. Coordinate System
Create CoordinateSystem SunBarycenter;
SunBarycenter.Origin = SolarSystemBarycenter;
SunBarycenter.Axes = MJ2000Eq;

% 2. Spacecraft Setup (Using Python Outputs)
Create Spacecraft satEJ; % Earth to Jupiter Leg
satEJ.CoordinateSystem = SunBarycenter;
satEJ.DisplayStateType = Cartesian;
satEJ.Epoch = '03 Nov 1976 01:41:50.000';
% -> PASTE PYTHON satEJ X,Y,Z,VX,VY,VZ HERE

Create Spacecraft satJS; % Jupiter to Saturn Leg
satJS.CoordinateSystem = SunBarycenter;
satJS.DisplayStateType = Cartesian;
satJS.Epoch = '06 Aug 1980 07:45:27.000';
% -> PASTE PYTHON satJS X,Y,Z,VX,VY,VZ HERE

% 3. Physics & Solvers
Create ForceModel SunProp_ForceModel;
SunProp_ForceModel.CentralBody = Sun;
SunProp_ForceModel.PointMasses = {Sun, Earth, Jupiter, Saturn};

Create Propagator SunProp;
SunProp.FM = SunProp_ForceModel;

Create DifferentialCorrector DefaultDC;
Create Yukon Yukon1;

Create Variable patchEpoch DeltaV_Penalty;
\end{lstlisting}

\newpage
% ==============================================================================
\section*{Step 2: The Targeter Implementation (Unoptimized)}

A Differential Corrector uses \texttt{Achieve} commands. Because we have six independent variables (three velocities for \texttt{satEJ} and three for \texttt{satJS}) but only three goals (match X, Y, and Z positions), the system is under-determined. The Targeter will arbitrarily change the velocities until the position gap closes, entirely ignoring the physics of a "free" gravity assist.

\begin{lstlisting}[caption=Mission Sequence: Targeter (DC)]
BeginMissionSequence;

patchEpoch = satJS.A1ModJulian; % The specific flyby date

Target DefaultDC {SolveMode = Solve, ExitMode = DiscardAndContinue};
    % Give the solver permission to tweak the departure and flyby vectors
    Vary DefaultDC(satEJ.VX = satEJ.VX, {Perturbation = 0.0001, MaxStep = 0.5});
    Vary DefaultDC(satEJ.VY = satEJ.VY, {Perturbation = 0.0001, MaxStep = 0.5});
    Vary DefaultDC(satEJ.VZ = satEJ.VZ, {Perturbation = 0.0001, MaxStep = 0.5});
    
    Vary DefaultDC(satJS.VX = satJS.VX, {Perturbation = 0.0001, MaxStep = 0.5});
    Vary DefaultDC(satJS.VY = satJS.VY, {Perturbation = 0.0001, MaxStep = 0.5});
    Vary DefaultDC(satJS.VZ = satJS.VZ, {Perturbation = 0.0001, MaxStep = 0.5});

    % Propagate Leg 1 to the Patch Point
    Propagate SunProp(satEJ) {satEJ.A1ModJulian = patchEpoch};

    % Force the positions to match
    Achieve DefaultDC(satEJ.SunBarycenter.X = satJS.SunBarycenter.X, {Tolerance = 0.1});
    Achieve DefaultDC(satEJ.SunBarycenter.Y = satJS.SunBarycenter.Y, {Tolerance = 0.1});
    Achieve DefaultDC(satEJ.SunBarycenter.Z = satJS.SunBarycenter.Z, {Tolerance = 0.1});
EndTarget;

% Post-Target Analysis: Calculate the Velocity Mismatch
DeltaV_Penalty = sqrt((satEJ.SunBarycenter.VX - satJS.SunBarycenter.VX)^2 + 
                      (satEJ.SunBarycenter.VY - satJS.SunBarycenter.VY)^2 + 
                      (satEJ.SunBarycenter.VZ - satJS.SunBarycenter.VZ)^2);
\end{lstlisting}

\begin{tcolorbox}[colback=red!5!white,colframe=red!75!black,title=Analysis Note]
If you run this sequence, check the final \texttt{DeltaV\_Penalty}. It will likely be hundreds of meters per second (or several kilometers per second). The Targeter mathematically closed the gap, but destroyed the energy efficiency in the process!
\end{tcolorbox}

% ==============================================================================
\section*{Step 3: The Yukon Optimization Implementation}

To fix the under-determined physics from Step 2, we swap the Targeter for the Yukon SQP Optimizer. We convert the \texttt{Achieve} commands into \texttt{NonLinearConstraints}, and we introduce a \texttt{Minimize} objective function. Now, GMAT will slide along the boundaries of the position constraint until it finds the absolute lowest $\Delta V$ combination.

\begin{lstlisting}[caption=Mission Sequence: Yukon Optimizer]
% (Replace the Target block with this Optimize block)

Optimize Yukon1 {SolveMode = Solve, ExitMode = SaveAndContinue};
    
    % The knobs remain the same
    Vary Yukon1(satEJ.VX = satEJ.VX, {Perturbation = 0.0001, MaxStep = 0.5});
    Vary Yukon1(satEJ.VY = satEJ.VY, {Perturbation = 0.0001, MaxStep = 0.5});
    Vary Yukon1(satEJ.VZ = satEJ.VZ, {Perturbation = 0.0001, MaxStep = 0.5});
    
    Vary Yukon1(satJS.VX = satJS.VX, {Perturbation = 0.0001, MaxStep = 0.5});
    Vary Yukon1(satJS.VY = satJS.VY, {Perturbation = 0.0001, MaxStep = 0.5});
    Vary Yukon1(satJS.VZ = satJS.VZ, {Perturbation = 0.0001, MaxStep = 0.5});

    Propagate SunProp(satEJ) {satEJ.A1ModJulian = patchEpoch};

    % 1. The Fences: Positions must match exactly
    NonLinearConstraint Yukon1(satEJ.SunBarycenter.X = satJS.SunBarycenter.X);
    NonLinearConstraint Yukon1(satEJ.SunBarycenter.Y = satJS.SunBarycenter.Y);
    NonLinearConstraint Yukon1(satEJ.SunBarycenter.Z = satJS.SunBarycenter.Z);

    % 2. The Objective Function: Calculate the mismatch magnitude
    DeltaV_Penalty = sqrt((satEJ.SunBarycenter.VX - satJS.SunBarycenter.VX)^2 + 
                          (satEJ.SunBarycenter.VY - satJS.SunBarycenter.VY)^2 + 
                          (satEJ.SunBarycenter.VZ - satJS.SunBarycenter.VZ)^2);

    % 3. The Goal: Drive the mismatch to zero (or as low as N-body physics allows)
    Minimize Yukon1(DeltaV_Penalty);

EndOptimize;

% If successful, propagate out to Saturn
Propagate SunProp(satJS) {satJS.ElapsedDays = 1095};
\end{lstlisting}

\end{document}